% =====================================================================
%  2. TRAVAUX CONNEXES
% =====================================================================
\section{Travaux connexes}\label{sec:connexes}

TRIADE emprunte à trois traditions~: la gestion de la dette technique, la
conception par contrat et, plus récemment, le génie logiciel assisté par IA
(AISE). Cette section les passe en revue successivement, puis positionne le
protocole par rapport à chacune.

\subsection{Dette technique}\label{sec:connexes-dette}

La métaphore de la dette technique, introduite par Cunningham, décrit la
conséquence d’un choix de conception opportun~: une facilité immédiate qui se
rembourse, avec intérêts, dans la maintenance future \cite{cunningham1992}.
Kruchten et al. en ont fait un objet de recherche à part entière, en
identifiant notamment la dimension \emph{sociale} et documentaire de la dette
— la perte de contexte des décisions passées \cite{kruchten2012}. La revue
systématique de Li et al. classe les activités de gestion en trois
familles~: identification, mesure et remboursement \cite{li2015}, et des
méthodes comme SQALE proposent d’évaluer le montant de la dette sur le code
existant \cite{letouzey2012}. Un courant connexe formalise la mémoire des
décisions elles-mêmes — les \emph{architecture decision records} de Nygard
\cite{nygard2011} — en réaction à la même perte de contexte.

Ces approches partagent un trait~: elles interviennent \emph{a posteriori},
sur un code dont la dette est déjà contractée. TRIADE inverse la perspective~:
la dette est traitée à la source, en rendant son principal explicite avant
l’écriture. La couche I est précisément l’artefact qui empêche la classe la
plus insidieuse de dette — le \emph{contrat non écrit} (valeur-seuil enfouie,
convention d’interface implicite, comportement dépendant de l’hôte) — d’être
contractée en silence~: ce qui n’est pas spécifié n’est pas codé, et ce qui
est codé renvoie à sa spécification. Le journal des décisions, requis par la
gouvernance du protocole, généralise les enregistrements de Nygard en les
rattachant à une fonction objectif explicite.

\subsection{Conception par contrat et spécification}\label{sec:connexes-dbc}

La lignée remonte à la logique axiomatique de Hoare, qui attache des
pré-conditions et post-conditions aux programmes \cite{hoare1969}, et à la
discipline de la programmation vérifiée de Dijkstra \cite{dijkstra1976}.
Meyer a popularisé la \emph{conception par contrat}~: chaque routine publie
ses obligations mutuelles — pré-conditions, post-conditions, invariants —
vérifiées à l’exécution \cite{meyer1992}~; Findler et Felleisen ont étendu la
notion aux fonctions d’ordre supérieur avec un contrôle dynamique des
violations \cite{findler2002}. Sur le plan de la vérification, QuickCheck a
établi que des propriétés déclaratives et des tests aléatoires contraints
formaient un oracle bon marché et redoutable \cite{claessen2000}. Pour le
code hérité, Feathers définit les \emph{characterization tests}~: figer le
comportement existant sous des tests avant d’y toucher \cite{feathers2004},
dans la tradition du remaniement contrôlé \cite{fowler1999}. Enfin, les
méthodes formelles légères (Alloy \cite{jackson2006}) ou lourdes (TLA+
\cite{lamport2002}) offrent des langages de spécification vérifiables par
model checking.

TRIADE se distingue des deux manières. D’une part, ses contrats sont
\emph{légers par décision explicite} (D3 du journal du projet)~: signatures
typées, pré/post-conditions en prose contrainte et invariants — pas de
sémantique formelle exécutable, au prix d’une vérification par tests plutôt
que par preuve. D’autre part et surtout, le contrat change de plan~: dans la
conception par contrat classique, le contrat est une annotation \emph{dans} le
code, surveillée à l’exécution~; dans TRIADE, le contrat est un artefact
\emph{extérieur} au code — diffable, versionné, antérieur à l’écriture — qui
\emph{conditionne le droit même de commiter} (règle R5). Le contrat n’est
plus une assertion runtime mais un instrument de gouvernance du processus.

\subsection{Génie logiciel assisté par IA}\label{sec:connexes-aise}

L’évaluation des modèles de code a d’abord mesuré des capacitudes~: réussite
sur des bancs d’essai comme HumanEval \cite{chen2021} ou la synthèse de
programmes à partir de spécifications en langage naturel \cite{austin2021}.
Un second courant mesure les effets sur les développeurs~: évaluations de
productivité sur des tâches contrôlées \cite{ziegler2022}, essai randomisé en
entreprise \cite{peng2023}, études d’utilisabilité mettant en évidence
l’écart entre attentes et expérience \cite{vaithilingam2022}, et analyses
fine-grained des modes d’interaction \cite{barke2023}. Les revues
systématiques agrègent ce corpus et y relèvent des angles morts récurrents~:
mesures de qualité au-delà des bancs d’essai, maintenance à long terme,
et confiance \cite{hou2024,fan2023}. Enfin, l’intégrité de la chaîne
d’approvisionnement a produit ses propres instruments~: constructions
reproductibles bit à bit \cite{lamb2021} et métadonnées de provenance signées
(in-toto) \cite{torresarias2019}.

À notre connaissance, la littérature AISE ne propose pas de protocole de
gouvernance du \emph{processus de développement} lui-même. Les travaux
existentants supposent l’outillage standard (répertoire versionné, revue de
code, intégration continue) et évaluent l’agent dans ce cadre hérité de
l’ère humaine. TRIADE suppose au contraire l’agent acquis et redessine le
cadre~: l’agent n’accède à l’artefact que par la carte (économie de
contexte), n’écrit qu’après vérification (R5), et voit ses évolutions
arbitrées par une fonction objectif explicite. Les verrouillages
cryptographiques s’inspirent des approches de provenance
\cite{torresarias2019,lamb2021}, mais protègent une relation inédite~: celle
entre l’\emph{index de développement} privé et l’artefact livré.

\subsection{Positionnement}\label{sec:connexes-position}

Le tableau~\ref{tab:position} résume le positionnement. Les approches de
dette technique mesurent un état et recommandent~; la conception par contrat
instrumente le code~; l’AISE évalue des modèles et des effets. TRIADE est le
seul des quatre à définir un \emph{protocole de processus}~: un ensemble
d’artefacts obligatoires, de règles à sanction mécanique et d’un ordre de
priorité non compensatoire, applicable à tout exécutant — humain ou agent.
Cette généralité est aussi sa limite~: rien, dans l’état actuel, ne prouve
que le protocole survive au contact d’équipes humaines nombreuses, dont les
incitations diffèrent de celles d’un agent (section~\ref{sec:menaces}).

\begin{table}[t]
\centering\small
\caption{Positionnement de TRIADE parmi les approches connexes.}
\label{tab:position}
\begin{tabularx}{\columnwidth}{@{}l Y Y@{}}
\toprule
Approche & Moment & Instrument pivot \\
\midrule
Dette technique \cite{li2015} & après coup & mesures, tableaux \\
Contrat (DbC) \cite{meyer1992} & à l’exécution & assertions dans le code \\
AISE \cite{hou2024} & évaluation & modèles, bancs d’essai \\
\textbf{TRIADE} (cet art.) & \textbf{avant écriture} & \textbf{specs + règles
à sanction} \\
\bottomrule
\end{tabularx}
\end{table}
